% Vietnamese translation for OI Public Library
% Source-PDF: 英文资料 ENGLISH MATERIALS/General Suffix Automaton Construction Algorithm and Space Bounds.pdf
\documentclass[11pt]{article}
\usepackage{oipl-vn}
\usepackage{tikz}
\usetikzlibrary{arrows.meta,positioning}
\renewcommand{\figurename}{Hình}

\newtheorem{definition}{Định nghĩa}
\newtheorem{lemma}{Bổ đề}
\newtheorem{proposition}{Mệnh đề}
\newtheorem{corollary}{Hệ quả}
\newenvironment{proof}{\paragraph{Chứng minh.}}{\hfill$\dashv$\par}

\newcommand{\eps}{\epsilon}
\newcommand{\End}{\operatorname{end\text{-}set}}
\newcommand{\suff}{\operatorname{suff}}
\newcommand{\undefinedsym}{\operatorname{undefined}}

\title{Thuật toán xây dựng ô-tô-mát hậu tố tổng quát và các chặn không gian}
\author{Mehryar Mohri, Pedro Moreno, Eugene Weinstein}
\date{}

\begin{document}
\maketitle

\origtitle{General Suffix Automaton Construction Algorithm and Space Bounds}
\sourcepath{english/general-suffix-automaton-construction-algorithm-and-space-bounds.pdf}

\begin{abstract}
Ô-tô-mát hậu tố và ô-tô-mát nhân tố là các cấu trúc dữ liệu hiệu quả để biểu
diễn chỉ mục đầy đủ của một tập xâu. Chúng là các ô-tô-mát hữu hạn đơn định
tối tiểu biểu diễn tập mọi hậu tố hoặc mọi xâu con của một tập xâu. Bài viết
này đưa ra một phân tích mới về kích thước của ô-tô-mát hậu tố và ô-tô-mát
nhân tố của một tập xâu. Cụ thể, ô-tô-mát hậu tố hoặc ô-tô-mát nhân tố của
tập xâu $U$ có nhiều nhất $2Q-2$ trạng thái, trong đó $Q$ là số nút của cây
tiền tố biểu diễn các xâu trong $U$. Chặn này cải thiện đáng kể so với chặn
$2\|U\|-1$ của Blumer và cộng sự (1987), trong đó $\|U\|$ là tổng độ dài các
xâu trong $U$. Tổng quát hơn, bài viết đưa ra các chặn mới cho kích thước của
ô-tô-mát hậu tố hoặc nhân tố của một ô-tô-mát theo kích thước của ô-tô-mát
ban đầu và độ dài lớn nhất của một hậu tố chung giữa các xâu mà nó chấp nhận.
Bài viết cũng mô tả chi tiết một thuật toán thời gian tuyến tính để xây dựng
ô-tô-mát hậu tố $S$ hoặc ô-tô-mát nhân tố $F$ của $U$ trong thời gian
$O(|S|)$. Thuật toán thật ra áp dụng được cho mọi ô-tô-mát đầu vào
hậu-tố-duy-nhất và tổng quát hóa chặt chẽ cách xây dựng trực tuyến chuẩn của
ô-tô-mát hậu tố cho một xâu đơn. Thuật toán cũng có thể dùng trực tiếp để tạo
oracle hậu tố hoặc oracle nhân tố của một tập xâu, vốn có nhiều tính chất hữu
ích trong khớp xâu. Phân tích cho thấy việc dùng ô-tô-mát nhân tố của ô-tô-mát
có thể khả thi trong các ứng dụng quy mô lớn; điều này cũng được củng cố bởi
thực nghiệm trên một bài toán nhận dạng nhạc với hơn $15{,}000$ bài hát.
\end{abstract}

\noindent\textbf{Từ khóa:} khớp xâu, khớp mẫu, chỉ mục, văn bản đảo,
ô-tô-mát hữu hạn, cây hậu tố, ô-tô-mát hậu tố, ô-tô-mát nhân tố, nhận dạng
nhạc.

\section{Giới thiệu}

Tìm kiếm mẫu trong những khối lượng rất lớn văn bản ngôn ngữ tự nhiên, chuỗi
sinh học và các chuỗi số hóa dễ truy cập khác là một vấn đề trung tâm của khoa
học máy tính. Bài toán có nhiều ứng dụng và đã được nghiên cứu rộng rãi
\cite{gusfield,crochemore-rytter}.

Bài viết này xét bài toán xây dựng một chỉ mục đầy đủ, hay tệp đảo, cho một
tập xâu được biểu diễn bằng ô-tô-mát hữu hạn. Khi số xâu lớn, chẳng hạn hàng
nghìn hoặc hàng triệu xâu, cả tập có thể được lưu gọn dưới dạng ô-tô-mát, đồng
thời cho phép cài đặt hiệu quả các thuật toán tìm kiếm và những thuật toán
khác \cite{allauzen-index,weinstein-music}. Trong nhiều ngữ cảnh như nhận dạng
tiếng nói hoặc trích xuất thông tin, toàn bộ tập xâu thường được cho trực tiếp
dưới dạng một ô-tô-mát.

Một cấu trúc dữ liệu hiệu quả và gọn để biểu diễn chỉ mục đầy đủ của một tập
xâu là \term{ô-tô-mát hậu tố}: ô-tô-mát đơn định tối tiểu biểu diễn tập mọi
hậu tố của tập xâu. Vì một xâu con là tiền tố của một hậu tố, ô-tô-mát hậu tố
có thể dùng để kiểm tra một xâu $x$ có phải là xâu con hay không trong thời
gian tuyến tính $O(|x|)$, là tối ưu. Tương tự cây hậu tố, ô-tô-mát hậu tố còn
có những tính chất thú vị khác trong các bài toán khớp xâu, khiến việc dùng và
xây dựng nó trở nên hấp dẫn \cite{gusfield,crochemore-rytter}. Một cấu trúc
tương tự để biểu diễn chỉ mục đầy đủ là \term{ô-tô-mát nhân tố}: ô-tô-mát đơn
định tối tiểu biểu diễn tập mọi nhân tố, tức mọi xâu con, của một tập xâu.
Ô-tô-mát nhân tố có cùng tính chất tìm kiếm tuyến tính tối ưu như ô-tô-mát hậu
tố, và không bao giờ lớn hơn nó.

Việc xây dựng và kích thước của ô-tô-mát nhân tố đã được phân tích riêng cho
trường hợp một xâu đơn \cite{blumer-smallest,crochemore-transducers}. Các tác
giả này chứng minh một kết quả đáng chú ý: kích thước ô-tô-mát nhân tố của
một xâu $x$ là tuyến tính; cụ thể, với $|x|>3$, nó có nhiều nhất $2|x|-2$
trạng thái và $3|x|-4$ cung. Họ cũng đưa ra thuật toán trực tuyến thời gian
tuyến tính để xây dựng ô-tô-mát nhân tố từ $x$. Những kết quả tương tự cũng
đúng cho ô-tô-mát hậu tố, tức ô-tô-mát đơn định tối tiểu chấp nhận đúng tập
hậu tố của một xâu.

Việc xây dựng và kích thước của ô-tô-mát nhân tố của một tập hữu hạn
$U=\{x_1,\ldots,x_m\}$ cũng đã được nghiên cứu \cite{blumer-complete}. Các tác
giả chỉ ra rằng có thể xây dựng một ô-tô-mát chấp nhận mọi nhân tố của $U$ với
nhiều nhất $2\|U\|-1$ trạng thái và $3\|U\|-3$ cung, trong đó
\[
  \|U\|=\sum_{i=1}^m |x_i|
\]
là tổng độ dài của mọi xâu trong $U$.

Bài viết này chứng minh một chặn tốt hơn đáng kể cho kích thước của ô-tô-mát
hậu tố hoặc ô-tô-mát nhân tố của một tập xâu. Cụ thể, chúng có nhiều nhất
$2Q-2$ trạng thái, trong đó $Q$ là số nút của cây tiền tố biểu diễn các xâu
trong $U$. Số nút $Q$ có thể nhỏ hơn rất nhiều so với $\|U\|$, nên chặn không
gian này cải thiện rõ rệt so với kết quả trước \cite{blumer-complete}. Tổng
quát hơn, chúng tôi đưa ra các chặn mới cho kích thước ô-tô-mát hậu tố hoặc
nhân tố của một ô-tô-mát hữu hạn phi chu trình, theo kích thước của ô-tô-mát
ban đầu và độ dài lớn nhất của một hậu tố chung giữa các xâu được chấp nhận.
Kết quả này có thể so sánh với kết quả của Inenaga và cộng sự cho đồ thị từ
phi chu trình có hướng dạng nén, có độ phức tạp $O(|\Sigma|Q)$ phụ thuộc vào
kích thước bảng chữ cái \cite{inenaga}.

Dựa trên phân tích chặn không gian, chúng tôi cũng đưa ra một thuật toán đơn
giản để xây dựng ô-tô-mát hậu tố $S$ hoặc ô-tô-mát nhân tố $F$ của $U$ trong
thời gian $O(|S|)$ từ cây tiền tố biểu diễn $U$. Thuật toán thật ra áp dụng cho
mọi ô-tô-mát đầu vào hậu-tố-duy-nhất và tổng quát hóa chặt chẽ phép xây dựng
trực tuyến chuẩn của ô-tô-mát hậu tố cho một xâu đơn.

Động cơ ban đầu của công trình là thiết kế một hệ thống nhận dạng nhạc quy mô
lớn \cite{weinstein-music,mohri-robust}. Trong hệ thống đó, cơ sở dữ liệu bài
hát được biểu diễn bằng một ô-tô-mát hữu hạn gọn. Để tìm kiếm nhanh các đoạn
nhạc ngắn, chúng tôi xây dựng ô-tô-mát nhân tố đơn định tối tiểu của ô-tô-mát
bài hát. Thực nghiệm cho thấy kích thước ô-tô-mát nhân tố không quá lớn. Tuy
nhiên, để bảo đảm khả năng mở rộng tới tập bài hát lớn hơn, chẳng hạn vài
triệu bài, ta cần một chặn cho kích thước ô-tô-mát nhân tố của ô-tô-mát. Trong
ứng dụng này cũng như trong nhiều ứng dụng khác, các xâu gốc không chia sẻ các
hậu tố dài. Điều đó dẫn tới phân tích kích thước ô-tô-mát nhân tố theo độ dài
các hậu tố chung trong ô-tô-mát ban đầu.

Phần còn lại của bài viết được tổ chức như sau. Mục 2 giới thiệu các định
nghĩa và thuật ngữ về xâu và ô-tô-mát. Mục 3 trình bày phân tích mới về
ô-tô-mát nhân tố và các chặn mới cho kích thước ô-tô-mát hậu tố và nhân tố của
một ô-tô-mát. Mục 4 mô tả chi tiết thuật toán thời gian tuyến tính để xây dựng
ô-tô-mát hậu tố và nhân tố của một tập hữu hạn các xâu, hoặc của một
ô-tô-mát hậu-tố-duy-nhất bất kỳ, kèm mã giả. Mục 5 mô tả ngắn việc dùng
ô-tô-mát nhân tố trong nhận dạng nhạc và báo cáo một số kết quả thực nghiệm
về kích thước của chúng.

\section{Nhân tố của một ô-tô-mát hữu hạn}

Mục này nhắc lại một số tính chất then chốt của các nhân tố của một ô-tô-mát
hữu hạn cố định, tổng quát hóa các quan sát tương tự của Blumer và cộng sự cho
một xâu đơn \cite{blumer-complete}.

Ký hiệu $\Sigma$ là một bảng chữ cái hữu hạn. Độ dài của xâu $x\in\Sigma^*$
trên bảng chữ cái đó được ký hiệu là $|x|$. Một \term{nhân tố}, hay xâu con,
của $x$ là một dãy ký hiệu xuất hiện liên tiếp trong $x$. Vì vậy, $y$ là nhân
tố của $x$ khi và chỉ khi tồn tại $u,v\in\Sigma^*$ sao cho $x=uyv$. Một
\term{hậu tố} của $x$ là một nhân tố xuất hiện ở cuối $x$; tương đương, $y$ là
hậu tố của $x$ khi và chỉ khi tồn tại $u\in\Sigma^*$ sao cho $x=uy$. Tương tự,
$y$ là \term{tiền tố} của $x$ khi và chỉ khi tồn tại $u\in\Sigma^*$ sao cho
$x=yu$. Tổng quát hơn, nhân tố, hậu tố hoặc tiền tố của một tập xâu $U$, hoặc
của một ô-tô-mát $A$, là nhân tố, hậu tố hoặc tiền tố của một xâu trong $U$,
hoặc của một xâu được $A$ chấp nhận. Ký hiệu $\eps$ biểu diễn xâu rỗng; với
mọi $x\in\Sigma^*$, $\eps$ luôn là tiền tố, hậu tố và nhân tố của $x$.

\begin{definition}
Cho $k$ là một số nguyên không âm. Ta nói ô-tô-mát hữu hạn $A$ là
\term{$k$-hậu-tố-duy-nhất} nếu không có hai xâu nào được $A$ chấp nhận cùng
chia sẻ một hậu tố độ dài $k$. Ta gọi $A$ là \term{hậu-tố-duy-nhất} khi nó là
$k$-hậu-tố-duy-nhất với $k=1$.
\end{definition}

Ví dụ chạy trong bài là một ô-tô-mát đơn giản chấp nhận ba xâu \texttt{ac},
\texttt{acab}, \texttt{acba}; ba xâu kết thúc bằng các ký hiệu khác nhau, nên
ô-tô-mát này là hậu-tố-duy-nhất.

Các kết quả chính của bài viết đúng cho ô-tô-mát hậu-tố-duy-nhất, nhưng một
số kết quả cũng được trình bày cho trường hợp tổng quát của ô-tô-mát phi chu
trình bất kỳ. Ký hiệu $F(A)$ là ô-tô-mát đơn định tối tiểu chấp nhận tập các
nhân tố của ô-tô-mát hữu hạn $A$, tức tập các nhân tố của những xâu được $A$
chấp nhận. Tương tự, $S(A)$ là ô-tô-mát đơn định tối tiểu chấp nhận tập các
hậu tố của $A$.

\begin{definition}
Cho $A$ là một ô-tô-mát hữu hạn. Với mọi $x\in\Sigma^*$, định nghĩa
$\End(x)$ là tập trạng thái của $A$ đạt được bởi những đường đi trong $A$ bắt
đầu bằng $x$. Ta nói hai xâu $x,y\in\Sigma^*$ tương đương, ký hiệu $x\equiv y$,
khi $\End(x)=\End(y)$. Đây là một quan hệ tương đương bất biến phải trên
$\Sigma^*$. Ký hiệu $[x]$ là lớp tương đương của $x$.
\end{definition}

\begin{lemma}
Giả sử $A$ là hậu-tố-duy-nhất. Khi đó, một nhân tố $x$ không phải hậu tố của
$A$ là phần tử dài nhất của $[x]$ khi và chỉ khi $x$ là tiền tố của $A$, hoặc
cả $ax$ và $bx$ đều là nhân tố của $A$ với hai ký hiệu phân biệt
$a,b\in\Sigma$.
\end{lemma}

\begin{proof}
Cho $x$ là một nhân tố không phải hậu tố của $A$. Nếu $x$ không phải tiền tố,
rõ ràng phải có hai ký hiệu phân biệt $a,b$ sao cho $ax$ và $bx$ là nhân tố
của $A$; nếu không, $[x]$ sẽ chứa một phần tử dài hơn.

Ngược lại, giả sử $ax$ và $bx$ đều là nhân tố của $A$ với $a\ne b$. Gọi $y$ là
phần tử dài nhất của $[x]$. Lấy $q\in\End(x)=\End(y)$. Vì $x$ không phải hậu
tố, $q$ không phải trạng thái kết thúc, nên tồn tại một xâu không rỗng $z$ gán
nhãn cho một đường đi từ $q$ tới trạng thái kết thúc. Do $A$ là
hậu-tố-duy-nhất, cả $xz$ và $yz$ là hậu tố của cùng một xâu. Vì $y$ dài nhất
trong $[x]$, $x$ phải là hậu tố của $y$. Nhưng do $ax$ và $bx$ cùng là nhân tố
với $a\ne b$, suy ra bắt buộc $y=x$. Cuối cùng, nếu $x$ là tiền tố, hiển nhiên
nó là phần tử dài nhất của $[x]$.
\end{proof}

Với ô-tô-mát ở ví dụ chạy trên, ô-tô-mát hậu tố $S(A)$ có thể đọc như sau:
các trạng thái của $S(A)$ chính là những lớp tương đương của các hậu tố và các
tiền tố của chúng.

\begin{proposition}
Giả sử $A$ là hậu-tố-duy-nhất. Xét ô-tô-mát đơn định
$S_A=(Q_S,I_S,F_S,E_S)$ có tập trạng thái
\[
  Q_S=\{[x]\ne\emptyset:x\in\Sigma^*\},
\]
trạng thái đầu $I_S=\{[\eps]\}$, tập trạng thái kết thúc
\[
  F_S=\{[x]:\End(x)\cap F_A\ne\emptyset\},
\]
trong đó $F_A$ là tập trạng thái kết thúc của $A$, và tập cung
\[
  E_S=\{([x],a,[xa]): [x],[xa]\in Q_S\}.
\]
Khi đó $S_A$ là ô-tô-mát hậu tố đơn định tối tiểu của $A$, tức $S_A=S(A)$.
\end{proposition}

\begin{proof}
Theo xây dựng, $S_A$ đơn định và chấp nhận đúng tập hậu tố của $A$. Cho $[x]$
và $[y]$ là hai trạng thái tương đương của $S_A$. Khi đó, với mọi
$z\in\Sigma^*$, $[xz]\in F_A$ khi và chỉ khi $[yz]\in F_A$, nghĩa là $xz$ là
hậu tố của $A$ khi và chỉ khi $yz$ là hậu tố của $A$. Vì $A$ là
hậu-tố-duy-nhất, điều này kéo theo $x$ là hậu tố của $y$ hoặc ngược lại; do đó
$[x]=[y]$. Vì vậy $S_A$ tối tiểu.
\end{proof}

Trong phần sau, ta quan tâm tới trường hợp ô-tô-mát $A$ là phi chu trình. Ký
hiệu $|A|_Q$ là số trạng thái của $A$, $|A|_E$ là số cung của $A$, và $|A|$ là
kích thước của $A$, tức tổng số trạng thái và số cung.

\section{Chặn không gian cho ô-tô-mát nhân tố}

Mục tiêu của mục này là suy ra các chặn mới cho kích thước của $S(A)$ và
$F(A)$ trong trường hợp quan trọng với ứng dụng của chúng tôi: $A$ là một
ô-tô-mát phi chu trình, thường đơn định và tối tiểu, biểu diễn một tập xâu.

Khi $A$ biểu diễn một xâu đơn, đã có các thuật toán chuẩn xây dựng $S(A)$ và
$F(A)$ từ $A$ trong thời gian tuyến tính \cite{blumer-smallest,crochemore-transducers}.
Trong trường hợp tổng quát, $S(A)$ có thể được xây dựng từ $A$ như sau: thêm
một cung $\eps$ từ trạng thái đầu của $A$ tới mỗi trạng thái của $A$, sau đó
khử $\eps$, đơn định hóa và tối tiểu hóa. $F(A)$ có thể thu được tương tự bằng
cách thêm bước biến mọi trạng thái thành kết thúc trước khi khử $\eps$, đơn
định hóa và tối tiểu hóa. Nó cũng có thể thu được từ $S(A)$ bằng cách biến mọi
trạng thái của $S(A)$ thành kết thúc rồi tối tiểu hóa.

Khi $A$ biểu diễn một xâu đơn $x$, kích thước của $S(A)$ và $F(A)$ là tuyến
tính theo $|x|$. Chính xác hơn \cite{crochemore-transducers,blumer-smallest},
\[
\begin{array}{ll}
  |S(A)|_Q \le 2|x|-1, & |S(A)|_E \le 3|x|-4,\\
  |F(A)|_Q \le 2|x|-2, & |F(A)|_E \le 3|x|-4.
\end{array}
\tag{1}
\]
Các chặn này là chặt với các xâu có độ dài lớn hơn ba. Blumer và cộng sự
\cite{blumer-complete} đưa ra kết quả tương tự cho tập xâu $U$:
\[
  |F(U)|_Q \le 2\|U\|-1,\qquad |F(U)|_E \le 3\|U\|-3.
\tag{2}
\]

Nhìn chung, kích thước của một ô-tô-mát phi chu trình $A$ biểu diễn một tập
hữu hạn $U$ có thể nhỏ hơn rất nhiều so với $\|U\|$; thật ra $|A|$ có thể nhỏ
hơn theo cấp số mũ. Vì thế, ta muốn chặn kích thước $S(A)$ hoặc $F(A)$ theo
kích thước của $A$, thay vì theo tổng độ dài mọi xâu được $A$ chấp nhận.

Với mỗi trạng thái $q$ của $S(A)$, ký hiệu $\suff(q)$ là tập các xâu gán nhãn
cho những đường đi từ $q$ tới một trạng thái kết thúc. Ký hiệu $N(q)$ là tập
các trạng thái trong $A$ mà từ đó có một đường đi gán nhãn bằng một xâu không
rỗng thuộc $\suff(q)$ tới một trạng thái kết thúc.

\begin{lemma}
Cho $A$ là ô-tô-mát hậu-tố-duy-nhất và $q,q'$ là hai trạng thái của $S(A)$ sao
cho $N(q)\cap N(q')\ne\emptyset$. Khi đó
\[
\begin{aligned}
  &\bigl(\suff(q)\subseteq\suff(q')\ \text{và}\ N(q)\subseteq N(q')\bigr)\\
  &\quad\text{hoặc}\quad
  \bigl(\suff(q')\subseteq\suff(q)\ \text{và}\ N(q')\subseteq N(q)\bigr).
\end{aligned}
\tag{3}
\]
\end{lemma}

\begin{proof}
Vì $S(A)$ tối tiểu, mọi trạng thái của nó đều truy cập được từ trạng thái đầu.
Gọi $u$ là nhãn của một đường đi từ trạng thái đầu $I$ của $S(A)$ tới $q$, và
tương tự gọi $u'$ là nhãn của một đường đi từ $I$ tới $q'$.

Theo giả thiết, tồn tại $p\in N(q)\cap N(q')$. Do đó tồn tại các xâu không
rỗng $v\in\suff(q)$ và $v'\in\suff(q')$ sao cho cả $v$ và $v'$ đều gán nhãn
cho đường đi từ $p$ tới trạng thái kết thúc.

Theo định nghĩa của $u$ và $u'$, cả $uv$ và $u'v'$ đều là hậu tố của $A$. Vì
$A$ là hậu-tố-duy-nhất và $v$ không rỗng, chỉ có một xâu được $A$ chấp nhận
kết thúc bằng $v$. Cũng chỉ có một xâu được $A$ chấp nhận kết thúc bằng $uv$.
Hai xâu đó phải trùng nhau.

Suy ra mọi xâu được $A$ chấp nhận và có hậu tố $v$ cũng có hậu tố $uv$. Đặc
biệt, nhãn của mọi đường đi từ trạng thái đầu tới $p$ phải có hậu tố $u$. Lập
luận tương tự với $v'$ cho thấy nhãn đó cũng phải có hậu tố $u'$. Vậy $u$ và
$u'$ là hậu tố của cùng một xâu; do đó $u$ là hậu tố của $u'$ hoặc ngược lại:
nếu $uv$ và $u'v$ là hậu tố của cùng một xâu $x$, thì $u$ và $u'$ là hai hậu
tố lồng nhau của $x$.

Không mất tính tổng quát, giả sử $u$ là hậu tố của $u'$. Khi đó với mọi $w$,
nếu $u'w$ là hậu tố của $A$ thì $uw$ cũng là hậu tố của $A$. Vì vậy
$\suff(q')\subseteq\suff(q)$, kéo theo $N(q')\subseteq N(q)$. Trường hợp
$u'$ là hậu tố của $u$ cho kết luận còn lại.
\end{proof}

Bổ đề 2 vẫn đúng ngay cả khi $A$ là ô-tô-mát không đơn định.

\begin{lemma}
Cho $A$ là ô-tô-mát đơn định hậu-tố-duy-nhất và $q,q'$ là hai trạng thái phân
biệt của $S(A)$ sao cho $N(q)=N(q')$. Khi đó hoặc $q$ là trạng thái kết thúc
còn $q'$ không phải, hoặc $q'$ là trạng thái kết thúc còn $q$ không phải.
\end{lemma}

\begin{proof}
Giả sử $N(q)=N(q')$. Theo Bổ đề 2, suy ra $\suff(q)=\suff(q')$. Do đó cùng các
xâu không rỗng gán nhãn cho các đường đi từ $q$ tới trạng thái kết thúc cũng
gán nhãn cho các đường đi từ $q'$ tới trạng thái kết thúc. Vì $S(A)$ tối tiểu,
hai trạng thái phân biệt $q$ và $q'$ không tương đương. Vậy một trong hai trạng
thái phải cho phép đường đi rỗng tới trạng thái kết thúc, còn trạng thái kia
thì không.
\end{proof}

\begin{proposition}
Cho $A$ là ô-tô-mát đơn định, tối tiểu, hậu-tố-duy-nhất, chấp nhận các xâu có
độ dài lớn hơn ba. Khi đó số trạng thái của ô-tô-mát hậu tố của $A$ thỏa
\[
  |S(A)|_Q \le 2|A|_Q-3.
\tag{4}
\]
\end{proposition}

\begin{proof}
Nếu mọi xâu được $A$ chấp nhận đều có dạng $a^n$, thì $S(A)$ thu được đơn giản
bằng cách biến mọi trạng thái của $A$ thành trạng thái kết thúc, và chặn hiển
nhiên đúng. Vì vậy, phần còn lại giả sử không phải mọi xâu được chấp nhận đều
có dạng này.

Gọi $F$ là trạng thái kết thúc duy nhất của $S(A)$ không có cung ra. Bổ đề 2
và 3 cho phép định nghĩa một cây $T$ trên mọi trạng thái của $S(A)$ trừ $F$,
theo thứ tự
\[
  N(q)\sqsubseteq N(q')
  \quad\Longleftrightarrow\quad
  \begin{cases}
    N(q)\subset N(q'),\quad\text{hoặc}\\
    N(q)=N(q')\ \text{và}\ q'\ \text{kết thúc, còn}\ q\ \text{không kết thúc.}
  \end{cases}
\tag{5}
\]
Ta đồng nhất mỗi nút của $T$ với trạng thái tương ứng trong $S(A)$. Theo Mệnh
đề 1, mỗi trạng thái $q$ của $S(A)$ cũng được đồng nhất với một lớp tương
đương $[x]$. Với $q\ne F$, nếu $[x]$ là lớp tương ứng thì, do $A$ là
hậu-tố-duy-nhất, $\End(x)$ trùng với $N(q)$.

Ta chứng minh số nút của $T$ nhiều nhất là $2|A|_Q-4$; từ đó suy ra chặn cần
chứng minh cho $S(A)$. Ta chặn riêng số nút phân nhánh và không phân nhánh
của $T$.

Cho $q$ là một nút của $T$, $[x]$ là lớp tương ứng, với $x$ là phần tử dài nhất
của lớp. Con của $q$ là các nút tương ứng với các lớp $[ax]$, trong đó
$a\in\Sigma$ và $ax$ là nhân tố của $A$.

Theo Bổ đề 1, nếu $x$ là nhân tố không phải hậu tố và không phải tiền tố, thì
tồn tại hai nhân tố $ax$ và $bx$ với $a\ne b$. Vì thế $q$ có ít nhất hai con
tương ứng với $[ax]$ và $[bx]$, nên là nút phân nhánh. Do đó nút không phân
nhánh chỉ có thể là nút mà $x$ là tiền tố, hoặc nút mà $x$ là hậu tố, tức trạng
thái kết thúc của $S(A)$.

Vì các xâu được $A$ chấp nhận không phải đều có dạng $a^n$, tiền tố rỗng
$\eps$ xuất hiện trong ít nhất hai ngữ cảnh trái khác nhau $a$ và $b$ với
$a\ne b$. Vì vậy tiền tố $\eps$, tương ứng với gốc của $T$, chắc chắn phân
nhánh. Gọi $f$ là trạng thái kết thúc duy nhất của $A$ không có cung ra. Lớp
tương đương của nhân tố dài nhất kết thúc tại $f$, tức xâu dài nhất được $A$
chấp nhận, tương ứng với trạng thái $F$ của $S(A)$ và không nằm trong $T$.
Vậy có nhiều nhất $|A|_Q-2$ tiền tố không phân nhánh.

Mặt khác, với mỗi xâu được $A$ chấp nhận có nhiều nhất một nút không phân
nhánh tương ứng với hậu tố. Gọi $N_{\rm str}$ là số xâu được $A$ chấp nhận.
Khi đó số nút không phân nhánh $N_{\rm nb}$ của $T$ thỏa
\[
  N_{\rm nb}\le |A|_Q-2+N_{\rm str}.
\]

Để chặn số nút phân nhánh $N_b$, quan sát rằng do $A$ là hậu-tố-duy-nhất, mỗi
xâu được $A$ chấp nhận kết thúc bằng một ký hiệu riêng $a_i$,
$i=1,\ldots,N_{\rm str}$. Mỗi $a_i$ là một ngữ cảnh trái khác nhau của nhân tố
rỗng $\eps$, nên nút gốc $[\eps]$ có các con $[a_i]$. Gọi $T_{a_i}$ là cây con
gốc $[a_i]$ và $n_{a_i}$ là số lá của nó. Gọi
$a_j$, $j=N_{\rm str}+1,\ldots,N_{\rm str}+k$, là các con khác của gốc, và
$T_{a_j}$ là các cây con tương ứng. Một cây có $n_{a_i}$ lá có ít hơn
$n_{a_i}$ nút phân nhánh, nên số nút phân nhánh của $T_{a_i}$ nhiều nhất là
$n_{a_i}-1$. Tổng số lá của $T$ nhiều nhất là số tập con rời nhau của $Q$, bỏ
trạng thái đầu và $f$.

Nếu gốc chỉ có các con $[a_i]$, tức $k=0$, thì tồn tại ít nhất một $a_i$, gọi
là $a_1$, cũng là tiền tố của $A$; nếu có ký hiệu khác thì nó đã là con của
gốc. Nút $a_1$ khi đó cũng có con vì nó tương ứng với một hậu tố, tức một
trạng thái kết thúc của $S(A)$. Do đó $a_1$ không thể là lá trong trường hợp
này. Vì vậy
\[
  \sum_{i=1}^{N_{\rm str}+k} n_{a_i}\le |A|_Q-2-\mathbf 1_{k=0}.
\]
Tổng số nút phân nhánh của $T$, kể cả gốc, nhiều nhất là
\[
  N_b \le \sum_{i=1}^{N_{\rm str}+k}(n_{a_i}-1)+1
      \le |A|_Q-2-\mathbf 1_{k=0}-(N_{\rm str}+k)+1
      \le |A|_Q-2-N_{\rm str}.
\]
Suy ra tổng số nút của $T$ nhiều nhất là
$N_{\rm nb}+N_b\le 2|A|_Q-4$.
\end{proof}

Trong trường hợp riêng $A$ biểu diễn một xâu đơn $x$, chặn của Mệnh đề 2 khớp
với chặn trong \cite{crochemore-transducers,blumer-smallest}, vì
$|A|_Q=|x|+1$. Chặn này là chặt với các xâu có độ dài lớn hơn ba, nên cũng
chặt cho các ô-tô-mát chấp nhận xâu có độ dài lớn hơn ba. Ô-tô-mát ở ví dụ
chạy là hậu-tố-duy-nhất, đơn định và tối tiểu, có $|A|_Q=6$ trạng thái; ô-tô
mát hậu tố tối tiểu của nó có $|S(A)|_Q=7<2|A|_Q-3$.

\begin{corollary}
Cho $A$ là ô-tô-mát đơn định, tối tiểu, hậu-tố-duy-nhất, chấp nhận các xâu có
độ dài lớn hơn ba. Khi đó số trạng thái của ô-tô-mát nhân tố của $A$ thỏa
\[
  |F(A)|_Q \le 2|A|_Q-3.
\tag{6}
\]
\end{corollary}

\begin{proof}
Như đã nói, $F(A)$ thu được từ $S(A)$ bằng cách biến mọi trạng thái thành kết
thúc rồi tối tiểu hóa. Do đó $|F(A)|\le |S(A)|$, và kết quả suy ra từ Mệnh đề
2.
\end{proof}

Blumer và cộng sự (1987) chứng minh rằng một ô-tô-mát chấp nhận mọi nhân tố
của một tập xâu $U$ có nhiều nhất $2\|U\|-1$ trạng thái, trong đó $\|U\|$ là
tổng độ dài các xâu trong $U$ \cite{blumer-complete}. Hệ quả sau cho một chặn
tốt hơn đáng kể theo số nút của cây tiền tố biểu diễn $U$, thường nhỏ hơn
nhiều so với $\|U\|$.

\begin{corollary}
Cho $U=\{x_1,\ldots,x_m\}$ là một tập xâu có độ dài lớn hơn ba và $A$ là cây
tiền tố biểu diễn $U$. Khi đó số trạng thái của ô-tô-mát nhân tố $F(U)$ và
ô-tô-mát hậu tố $S(U)$ của các xâu trong $U$ thỏa
\[
  |F(U)|_Q \le 2|A|_Q-2,\qquad
  |S(U)|_Q \le 2|A|_Q-2.
\tag{7}
\]
\end{corollary}

\begin{proof}
Gọi $B$ là cây tiền tố biểu diễn tập
$U'=\{x_1\mathtt{\$}_1,\ldots,x_m\mathtt{\$}_m\}$, thu được bằng cách gắn vào
mỗi xâu của $U$ một ký hiệu mới $\mathtt{\$}_i$, $i=1,\ldots,m$, để làm các
hậu tố phân biệt; gọi $B'$ là ô-tô-mát thu được bằng cách tối tiểu hóa $B$.
Theo xây dựng, $B$ có nhiều hơn $A$ đúng $m$ trạng thái, nhưng vì mọi trạng
thái kết thúc của $B$ tương đương và bị nhập lại sau khi tối tiểu hóa, $B'$ có
nhiều nhất hơn $A$ một trạng thái.

Theo xây dựng, $B'$ là ô-tô-mát hậu-tố-duy-nhất. Theo Mệnh đề 2,
$|S(B')|_Q\le 2|B'|_Q-3$. Xóa khỏi $S(B')$ các cung gán nhãn bởi các ký hiệu
phụ $\mathtt{\$}_i$ rồi nối ô-tô-mát kết quả ta thu được ô-tô-mát hậu tố tối
tiểu $S(U)$. Trong $S(B')$, phải có một trạng thái kết thúc chỉ đạt được qua
các cung gán nhãn $\mathtt{\$}_i$; trạng thái này trở nên không truy cập được
sau khi xóa ký hiệu phụ. Vì thế $S(U)$ có ít nhất một trạng thái ít hơn
$S(B')$, và
\[
  |S(U)|_Q \le |S(B')|_Q-1 \le 2|B'|_Q-4 = 2|A|_Q-2.
\tag{8}
\]
Chặn tương tự cho $F(U)$ suy ra theo lập luận trong chứng minh Hệ quả 1.
\end{proof}

Khi $A$ là $k$-hậu-tố-duy-nhất với $k$ tương đối nhỏ, như trong các ứng dụng
chúng tôi quan tâm, mệnh đề sau cho một chặn tiện dụng cho kích thước
ô-tô-mát hậu tố.

\begin{proposition}
Cho $A$ là ô-tô-mát đơn định $k$-hậu-tố-duy-nhất chấp nhận các xâu có độ dài
lớn hơn ba, và $n$ là số xâu được $A$ chấp nhận. Khi đó số trạng thái của
ô-tô-mát hậu tố của $A$ thỏa
\[
  |S(A)|_Q \le 2|A_k|_Q + 2kn - 3,
\tag{9}
\]
trong đó $A_k$ là phần của ô-tô-mát $A$ thu được bằng cách bỏ các trạng thái
và cung của mọi hậu tố độ dài $k$.
\end{proposition}

\begin{proof}
Cho $A$ là ô-tô-mát đơn định $k$-hậu-tố-duy-nhất chấp nhận các xâu có độ dài
lớn hơn ba, và mở rộng bảng chữ cái $\Sigma$ bằng $n$ ký hiệu tạm
$\mathtt{\$}_1,\ldots,\mathtt{\$}_n$. Bằng cách đánh dấu mỗi xâu được $A$ chấp
nhận bằng một ký hiệu riêng $\mathtt{\$}_i$, ta biến $A$ thành một ô-tô-mát đơn
định hậu-tố-duy-nhất $A'$.

Để làm vậy, trước hết ta trải phẳng mọi hậu tố độ dài $k$ của $A$. Trong
trường hợp xấu nhất, mọi hậu tố phân biệt này trước đó cùng chia sẻ một hậu tố
độ dài $k-1$; việc trải phẳng khi đó có thể làm tăng số trạng thái của $A$
thêm tối đa $kn-n$ trạng thái. Đánh dấu cuối mỗi hậu tố bằng một ký hiệu
$\mathtt{\$}$ riêng làm tăng thêm $n$ trạng thái. Ô-tô-mát kết quả $A'$ là đơn
định và $|A'|_Q\le |A_k|_Q+kn$. Theo Mệnh đề 2,
$|S(A')|_Q\le 2|A'|_Q-3$. Vì các cung gán nhãn $\mathtt{\$}$ chỉ có thể xuất
hiện ở cuối đường đi thành công trong $S(A')$, ta có thể xóa các cung này,
biến trạng thái xuất phát của chúng thành kết thúc, rồi tối tiểu hóa ô-tô-mát
thu được để nhận ô-tô-mát đơn định $A''$ chấp nhận tập hậu tố của $A$. Kết
luận của mệnh đề suy ra từ $|A''|_Q\le |S(A')|_Q$.
\end{proof}

Vì kích thước $F(A)$ luôn không lớn hơn kích thước $S(A)$, ta có ngay hệ quả
sau.

\begin{corollary}
Cho $A$ là ô-tô-mát $k$-hậu-tố-duy-nhất chấp nhận các xâu có độ dài lớn hơn
ba. Khi đó ô-tô-mát nhân tố của $A$ thỏa
\[
  |F(A)|_Q \le 2|A_k|_Q + 2kn - 3.
\tag{10}
\]
\end{corollary}

Chặn trong hệ quả này không chặt với các giá trị $k$ tương đối nhỏ: trong thực
tế, kích thước ô-tô-mát nhân tố không phụ thuộc vào $kn$, tổng độ dài các hậu
tố độ dài $k$, mà phụ thuộc nhiều hơn vào số trạng thái của $A$ dùng để biểu
diễn chúng; với ô-tô-mát tối tiểu, số này có thể nhỏ hơn đáng kể. Tuy nhiên,
với $k$ lớn, chẳng hạn khi mọi xâu có cùng độ dài và $k$ bằng đúng độ dài các
xâu được $A$ chấp nhận, chặn của chúng tôi trùng với chặn của
\cite{blumer-complete}.

\section{Thuật toán xây dựng ô-tô-mát hậu tố}

Mục này mô tả một thuật toán thời gian tuyến tính để xây dựng ô-tô-mát hậu tố
$S(A)$ của một ô-tô-mát đầu vào hậu-tố-duy-nhất $A$, hoặc tương tự là
ô-tô-mát nhân tố $F(A)$ của $A$. Vì ô-tô-mát nhân tố có thể thu được từ
$S(A)$ bằng cách biến mọi trạng thái của $S(A)$ thành kết thúc rồi áp dụng
thuật toán tối tiểu hóa ô-tô-mát đơn định phi chu trình thời gian tuyến tính
\cite{revuz}, chỉ cần mô tả thuật toán thời gian tuyến tính cho $S(A)$. Tuy
nhiên, cũng có thể đưa ra một thuật toán trực tiếp tương tự cho $F(A)$.

Các thuật toán dưới đây xây dựng $S(A)=(Q_S,I,F_S,\delta_S)$ từ
$A=(Q_A,I,F_A,\delta_A)$, trong đó $A$ là hậu-tố-duy-nhất,
$\delta_S:Q_S\times\Sigma\to Q_S$ là hàm chuyển bộ phận của $S(A)$, và
$\delta_A:Q_A\times\Sigma\to Q_A$ là hàm chuyển bộ phận của $A$. Như ở mục
trước, $f$ là trạng thái kết thúc của $A$ không có cung ra. Ta gọi đích của
liên kết hậu tố là \term{con trỏ hậu tố}.

\begin{center}
\begin{tabular}{r p{0.78\linewidth}}
\multicolumn{2}{l}{\textsc{Create-Suffix-Automaton}$(A,f)$}\\
1 & $S\leftarrow Q_S\leftarrow\{I\}$ \\
2 & $s[I]\leftarrow\undefinedsym$; $l[I]\leftarrow 0$\\
3 & \textbf{while} $S\ne\emptyset$ \textbf{do}\\
4 & \quad $p\leftarrow \operatorname{Head}(S)$\\
5 & \quad \textbf{for each} cung $(p,a,q)$ của $A$ \textbf{do}\\
6 & \quad\quad \textbf{if} $q\ne f$ \textbf{then}\\
7 & \quad\quad\quad $Q_S\leftarrow Q_S\cup\{q\}$\\
8 & \quad\quad\quad $l[q]\leftarrow l[p]+1$\\
9 & \quad\quad\quad \textsc{Suffix-Next}$(p,a,q)$\\
10 & \quad\quad\quad $\operatorname{Enqueue}(S,q)$\\
11 & $Q_S\leftarrow Q_S\cup\{f\}$\\
12 & \textbf{for each} trạng thái $p\in Q_A$ và $a\in\Sigma$ sao cho
      $\delta_A(p,a)=f$ \textbf{do}\\
13 & \quad \textsc{Suffix-Next}$(p,a,f)$\\
14 & \quad \textsc{Suffix-Final}$(f)$\\
15 & \textbf{for each} $p\in F_A$ \textbf{do}\\
16 & \quad \textsc{Suffix-Final}$(p)$\\
17 & \textbf{return} $S(A)=(Q_S,I,F_S,\delta_S)$
\end{tabular}
\end{center}

Thuật toán là một tổng quát hóa của phép xây dựng chuẩn cho một xâu đầu vào
sang trường hợp đầu vào là ô-tô-mát hậu-tố-duy-nhất. Nó duy trì hai giá trị
$s[q]$ và $l[q]$ cho mỗi trạng thái $q$ của $S(A)$. $s[q]$ là con trỏ hậu tố,
hay trạng thái thất bại, của $q$. $l[q]$ là độ dài đường đi dài nhất từ trạng
thái đầu tới $q$ trong $S(A)$. Giá trị $l$ được dùng để xác định các cung đặc,
hay chuyển tiếp đặc, trong quá trình xây dựng: cung $(p,a,q)$ là đặc nếu
$l[p]+1=l[q]$, tức nó nằm trên một đường đi dài nhất từ trạng thái đầu tới
$q$; nếu không, nó là cung tắt.

$S$ là hàng đợi chứa các trạng thái cần xét. Kỷ luật hàng đợi cụ thể không ảnh
hưởng tới tính đúng đắn; có thể giả sử dùng FIFO, tương ứng với duyệt theo
chiều rộng và có cài đặt tuyến tính. Trong mỗi vòng lặp dòng 3--10, một trạng
thái mới $p$ được lấy ra từ $S$. Việc xử lý các cung $(p,a,f)$ có đích là $f$
được hoãn tới giai đoạn sau (dòng 12--14). Lý do là $f$ có tính chất đặc biệt:
như đã thảo luận ở mục trước, nó có thể được xem là con của nhiều nút khác
nhau của cây $T$, nên có thể có các liên kết hậu tố khác nhau. Các cung khác
$(p,a,q)$ được xử lý lần lượt bằng cách tạo trạng thái đích $q$ nếu cần, thêm
nó vào $Q_S$, xác định $l[q]$ và gọi \textsc{Suffix-Next}$(p,a,q)$.

\begin{center}
\begin{tabular}{r p{0.78\linewidth}}
\multicolumn{2}{l}{\textsc{Suffix-Next}$(p,a,q)$}\\
1 & $l[q]\leftarrow\max(l[p]+1,l[q])$\\
2 & \textbf{while} $p\ne I$ và $\delta_S(p,a)=\undefinedsym$ \textbf{do}\\
3 & \quad $\delta_S(p,a)\leftarrow q$\\
4 & \quad $p\leftarrow s[p]$\\
5 & \textbf{if} $\delta_S(p,a)=\undefinedsym$ \textbf{then}\\
6 & \quad $\delta_S(I,a)\leftarrow q$\\
7 & \quad $s[q]\leftarrow I$\\
8 & \textbf{elseif} $l[p]+1=l[\delta_S(p,a)]$ và $\delta_S(p,a)\ne q$
    \textbf{then}\\
9 & \quad $s[q]\leftarrow \delta_S(p,a)$\\
10 & \textbf{else} $r\leftarrow q$\\
11 & \quad \textbf{if} $\delta_S(p,a)\ne q$ \textbf{then}\\
12 & \quad\quad $r\leftarrow$ bản sao của $\delta_S(p,a)$ với cùng các cung ra\\
13 & \quad\quad $Q_S\leftarrow Q_S\cup\{r\}$\\
14 & \quad\quad $s[q]\leftarrow r$\\
15 & \quad $s[r]\leftarrow s[\delta_S(p,a)]$\\
16 & \quad $s[\delta_S(p,a)]\leftarrow r$\\
17 & \quad $l[r]\leftarrow l[p]+1$\\
18 & \quad \textbf{while} $p\ne\undefinedsym$ và $l[\delta_S(p,a)]\ge l[r]$
     \textbf{do}\\
19 & \quad\quad $\delta_S(p,a)\leftarrow r$\\
20 & \quad\quad $p\leftarrow s[p]$
\end{tabular}
\end{center}

\begin{center}
\begin{tabular}{r p{0.78\linewidth}}
\multicolumn{2}{l}{\textsc{Suffix-Final}$(p)$}\\
1 & \textbf{if} $p\in F_S$ \textbf{then}\\
2 & \quad $p\leftarrow s[p]$\\
3 & \textbf{while} $p\ne\undefinedsym$ và $p\notin F_S$ \textbf{do}\\
4 & \quad $F_S\leftarrow F_S\cup\{p\}$\\
5 & \quad $p\leftarrow s[p]$
\end{tabular}
\end{center}

Thủ tục \textsc{Suffix-Next} xử lý mỗi cung $(p,a,q)$ rất giống phép xây dựng
ô-tô-mát hậu tố chuẩn cho một xâu. Vòng lặp dòng 2--4 duyệt các con trỏ hậu
tố lặp của $p$ chưa có cung ra gán nhãn $a$. Nó tạo các cung như vậy tới $q$
từ mọi con trỏ hậu tố lặp cho tới khi gặp trạng thái đầu hoặc một trạng thái
$p'$ đã có cung này. Trong trường hợp thứ nhất, con trỏ hậu tố của $q$ được
đặt là trạng thái đầu $I$ và cung $(I,a,q)$ được tạo.

Trong trường hợp thứ hai, nếu cung sẵn có $(p',a,q')$ là cung đặc và $q'=q$,
thì con trỏ hậu tố của $q$ được đặt là $q'$ (dòng 9). Nếu $q'\ne q$, tạo một
bản sao $r$ của $q'$ với cùng các cung ra (dòng 12) và đặt con trỏ hậu tố của
$q$ là $r$. Con trỏ hậu tố của $r$ được đặt thành $s[q']$ (dòng 15), con trỏ
của $q'$ được đặt thành $r$ (dòng 16), và $l[r]$ được định nghĩa là $l[p]+1$
(dòng 17). Các cung gán nhãn $a$ đi ra từ các con trỏ hậu tố lặp của $p$ được
xét và chuyển hướng tới $r$ miễn là chúng là các cung không đặc (dòng 18--20).

Thủ tục \textsc{Suffix-Final} đặt tính kết thúc của các trạng thái trong
$S(A)$. Với mọi trạng thái $p$ kết thúc trong $A$, trạng thái $p$ và mọi trạng
thái tìm được bằng cách đi theo chuỗi con trỏ hậu tố bắt đầu từ $p$ đều được
đưa vào $F_S$ trong vòng lặp dòng 3--5.

Có thể theo dõi thuật toán bằng cách ghi mỗi trạng thái trung gian dưới dạng
$(n/s,l)$, trong đó $n$ là số hiệu trạng thái, $s$ là con trỏ hậu tố của nó,
và $l$ là giá trị $l[n]$; các giai đoạn xây dựng $S(A)$ chỉ thay đổi các cung,
con trỏ hậu tố và giá trị $l$ theo đúng các dòng trong hai thủ tục trên.

Trong phép xây dựng \term{suffix oracle} \cite{allauzen-oracle}, không có
trạng thái mới nào được tạo so với đầu vào. Do đó oracle hậu tố của $A$ có thể
được xây dựng tương tự bằng cách thay dòng 12 trong \textsc{Suffix-Next} bởi
$r\leftarrow\delta_S(p,a)$ và bỏ các dòng 15--17. Thuật toán này vì thế mở
rộng trực tiếp phép xây dựng oracle hậu tố sang trường hợp đầu vào là
ô-tô-mát hậu-tố-duy-nhất.

Với kết quả độ phức tạp sau, ta giả sử hàm chuyển được biểu diễn hiệu quả để
có thể tìm trong thời gian $O(1)$ một cung ra có nhãn cụ thể tại mỗi trạng
thái. Một số tác giả giả sử biểu diễn danh sách kề và tìm kiếm nhị phân để tìm
cung tại một trạng thái, tốn $O(\min\{\log|\Sigma|,e_{\max}\})$, trong đó
$e_{\max}$ là bậc ra lớn nhất \cite{crochemore-transducers,crochemore-rytter}.
Nếu dùng giả thiết đó, các kết quả độ phức tạp ở đây cũng như của Blumer và
cộng sự \cite{blumer-smallest,blumer-complete} cần nhân thêm hệ số
$\min\{\log|\Sigma|,e_{\max}\}$.

Ta gọi việc chuyển hướng một cung đã từng được chuyển hướng trước đó là
\term{chuyển hướng nhiều lần}.

\begin{proposition}
Cho $A$ là ô-tô-mát đơn định tối tiểu hậu-tố-duy-nhất. Khi đó độ phức tạp thời
gian của thuật toán \textsc{Create-Suffix-Automaton}$(A,f)$ là $O(|S(A)|)$.
\end{proposition}

\begin{proof}
Ta phác thảo chứng minh. \textsc{Suffix-Next} được gọi nhiều nhất một lần cho
mỗi cung, nên tổng số lời gọi là $O(|A|)$. Cố định một cung $(p,a,q)$ của $A$
với $q\ne f$. Chi phí thực hiện các bước 1--20 của \textsc{Suffix-Next} tỉ lệ
với tổng số lần đi theo liên kết hậu tố lặp trong các vòng lặp dòng 2--4 và
18--20. Mỗi lần lặp của dòng 2--4 tạo ra một cung mới trong $S(A)$, nên tổng
số lần lặp này trên mọi lời gọi \textsc{Suffix-Next} là $O(|S(A)|)$.

Phân tích tổng số lần chuyển hướng trong vòng lặp dòng 18--20 dựa trên một mở
rộng của phân tích cho trường hợp một xâu đơn \cite{blumer-thesis,blumer-complete}.
Chặn tuyến tính cho tổng số lần chuyển hướng trong trường hợp một xâu đơn áp
dụng được cho trường hợp ô-tô-mát của chúng ta trên một chuỗi tuyến tính các
trạng thái của $A$. Từ tổ hợp xâu con cần thiết trong $A$ để gây ra chuyển
hướng, có thể chứng minh tổng số chuyển hướng nhiều lần là $O(|A|)$. Do đó độ
phức tạp tổng là $O(|S(A)|)$.
\end{proof}

\section{Ô-tô-mát nhân tố cho nhận dạng nhạc}

Chúng tôi kiểm chứng các nhận xét trên về ô-tô-mát nhân tố trong ngữ cảnh một
hệ thống nhận dạng nhạc \cite{weinstein-music,mohri-robust}. Nhận dạng nhạc là
nhiệm vụ khớp một luồng âm thanh với một bài hát cụ thể. Trong hệ thống này,
ta học một kho đơn vị \term{music phone} tương tự âm vị trong tiếng nói, cùng
một dãy music phone duy nhất đặc trưng cho mỗi bài hát. Tập music phone được
xem như bảng chữ cái, còn các dãy music phone là tập xâu; bài toán trở thành
một bài toán nhận dạng nhân tố. Cách tiếp cận là xây dựng một bộ chuyển đổi
gọn ánh xạ các dãy music phone tới định danh bài hát tương ứng.

\subsection{Xây dựng bộ chuyển đổi nhân tố}

Ký hiệu $\Sigma$ là tập music phone và tập dãy music phone mô tả $m$ bài hát
là $U=\{x_1,\ldots,x_m\}$, với $x_i\in\Sigma^*$ cho
$i\in\{1,\ldots,m\}$. Trong thực nghiệm, $m=15{,}455$, $|\Sigma|=1{,}024$ và
độ dài trung bình của một phiên âm $x_i$ lớn hơn $1{,}700$. Vì vậy trong
trường hợp xấu nhất có thể có tới
\[
  15{,}455\times 1{,}700^2 \approx 45\times 10^9
\]
nhân tố. Một biểu diễn ngây thơ dựa trên cây tiền tố sẽ quá lớn, nên ta biểu
diễn tập nhân tố bằng ô-tô-mát nhân tố gọn hơn nhiều. Chúng tôi xây dựng một
ô-tô-mát đơn định tối tiểu biểu diễn các dãy trong $U$, rồi xây dựng một bộ
chuyển đổi hữu hạn trạng thái đơn định tối tiểu ánh xạ mỗi bài hát tới định
danh của nó bằng các thuật toán đơn định hóa và tối tiểu hóa bộ chuyển đổi
\cite{mohri-transducers,mohri-applied}.

Gọi $T_0$ là bộ chuyển đổi chưa tối ưu ánh xạ các dãy phone tới định danh bài
hát. Khi $U$ chỉ gồm vài bài hát ngắn, $T_0$ có thể hình dung như một cây tiền
tố có nhãn đầu ra ở cuối mỗi bài hát:

\begin{figure}[ht]
\centering
\begin{tikzpicture}[
  >=Latex,
  state/.style={circle,draw,inner sep=1.8pt,minimum size=6mm,font=\small},
  final/.style={circle,draw,double,inner sep=1.8pt,minimum size=6mm,font=\small},
  edge/.style={-{Latex[length=1.8mm]},thick,font=\small}
]
  \node[state] (q0) at (0,0) {$0$};
  \node[state] (q1) at (1.4,1.1) {};
  \node[state] (q2) at (2.8,1.1) {};
  \node[final] (q3) at (4.2,1.1) {};
  \node[state] (q4) at (1.4,-0.1) {};
  \node[final] (q5) at (2.8,-0.1) {};
  \node[state] (q6) at (1.4,-1.3) {};
  \node[state] (q7) at (2.8,-1.3) {};
  \node[final] (q8) at (4.2,-1.3) {};

  \draw[edge] (q0) -- node[above left] {$a:\eps$} (q1);
  \draw[edge] (q1) -- node[above] {$b:\eps$} (q2);
  \draw[edge] (q2) -- node[above] {$c:1$} (q3);
  \draw[edge] (q0) -- node[above] {$a:\eps$} (q4);
  \draw[edge] (q4) -- node[above] {$d:2$} (q5);
  \draw[edge] (q0) -- node[below left] {$e:\eps$} (q6);
  \draw[edge] (q6) -- node[below] {$b:\eps$} (q7);
  \draw[edge] (q7) -- node[below] {$c:3$} (q8);

  \node[draw=none,align=left,font=\small] at (5.6,0)
    {nhãn vào: music phone\\nhãn ra: định danh bài hát};
\end{tikzpicture}
\caption{Sơ đồ minh họa \(T_0\): mỗi đường từ gốc tới trạng thái kết thúc đọc một bài hát ngắn và xuất định danh của bài hát đó ở cuối đường.}
\end{figure}

Gọi $A$ là acceptor thu được bằng cách bỏ nhãn đầu ra của $T_0$. Ô-tô-mát nhân
tố gọn $F(A)$ được xây dựng như ở Mục 3: tạo các cung $\eps$ từ trạng thái đầu
của $A$ tới mọi trạng thái khác, biến mọi trạng thái thành kết thúc, rồi khử
$\eps$, đơn định hóa và tối tiểu hóa. Chú ý rằng $F(A)$ không xuất ra định
danh bài hát gắn với mỗi nhân tố.

Để mô tả bước tiếp theo, ta nhắc ngắn một số tính chất của ô-tô-mát có trọng
số. Một ô-tô-mát có trọng số được định nghĩa trên một nửa vành
$(K,\oplus,\otimes,\bar 0,\bar 1)$, quy định tập trọng số và các phép toán đại
số để kết hợp trọng số dọc theo một đường đi và giữa các đường đi. Nửa vành
nhiệt đới
\[
  (\mathbb R\cup\{-\infty,+\infty\},\min,+,+\infty,0)
\]
được dùng rộng rãi trong các lĩnh vực như xử lý tiếng nói và văn bản. Trong
nửa vành nhiệt đới, tổng trọng số mà ô-tô-mát gán cho một xâu $s$ là đường đi
có trọng số nhỏ nhất trong ô-tô-mát mang nhãn $s$, trong đó trọng số của một
đường đi được tính bằng tổng trọng số các cung trên đường đi đó.

Để xây dựng ô-tô-mát nhân tố vẫn giữ được định danh bài hát, ta tạo một
acceptor có trọng số gọn trên nửa vành nhiệt đới, chấp nhận các nhân tố của
$U$ và gán tổng trọng số $s_x$ cho mỗi nhân tố $x$. Ưu điểm quan trọng của biểu
diễn này là có thể dùng đơn định hóa và tối tiểu hóa có trọng số
\cite{mohri-transducers}; trong các phép toán đó, định danh bài hát được xem
như trọng số có thể phân bố dọc đường đi. Các phép toán bảo toàn tính chất rằng
tổng trọng số trên đường đi mang nhãn $x$ là $s_x$. Gọi $F_w(A)$ là ô-tô-mát
được xây dựng tương tự $F(A)$, nhưng mỗi cung $\eps$ được thêm vào có trọng số
là định danh bài hát tương ứng. Sau khi đơn định hóa và tối tiểu hóa trên nửa
vành nhiệt đới, acceptor có trọng số $F_w(A)$ được biến thành bộ chuyển đổi
nhận dạng bài hát $T$ bằng cách xem mỗi trọng số đầu ra nguyên là một ký hiệu
đầu ra. Với một dãy music phone đầu vào, định danh bài hát tương ứng được thu
bằng cách cộng các đầu ra do $T$ sinh ra.

\subsection{Kích thước ô-tô-mát}

Sơ đồ dưới đây tóm tắt khác biệt giữa acceptor nhân tố không trọng số $F(A)$
và acceptor nhân tố có trọng số $F_w(A)$ cho hai bài hát. Trong bản có trọng
số, định danh bài hát được đưa vào trọng số trên các cung $\eps$ ban đầu, nên
thông tin bài hát vẫn được giữ sau đơn định hóa và tối tiểu hóa mà không làm
tăng kích thước cấu trúc một cách đáng kể.

\begin{figure}[ht]
\centering
\begin{tikzpicture}[
  >=Latex,
  state/.style={circle,draw,inner sep=1.8pt,minimum size=6mm,font=\small},
  edge/.style={-{Latex[length=1.8mm]},thick,font=\small}
]
  \node[draw=none] at (1.9,1.7) {\(F(A)\): chỉ nhận nhân tố};
  \node[state] (a0) at (0,0) {$0$};
  \node[state] (a1) at (1.6,0.8) {};
  \node[state] (a2) at (3.2,0.8) {};
  \node[state] (a3) at (1.6,-0.8) {};
  \node[state] (a4) at (3.2,-0.8) {};
  \draw[edge] (a0) -- node[above] {$b$} (a1);
  \draw[edge] (a1) -- node[above] {$c$} (a2);
  \draw[edge] (a0) -- node[below] {$d$} (a3);
  \draw[edge] (a3) -- node[below] {$c$} (a4);
  \draw[edge,bend left=18] (a1) to node[right] {$a$} (a4);

  \begin{scope}[xshift=6.1cm]
    \node[draw=none] at (1.9,1.7) {\(F_w(A)\): nhân tố kèm trọng số};
    \node[state] (w0) at (0,0) {$0$};
    \node[state] (w1) at (1.6,0.8) {};
    \node[state] (w2) at (3.2,0.8) {};
    \node[state] (w3) at (1.6,-0.8) {};
    \node[state] (w4) at (3.2,-0.8) {};
    \draw[edge] (w0) -- node[above] {$b/1$} (w1);
    \draw[edge] (w1) -- node[above] {$c/0$} (w2);
    \draw[edge] (w0) -- node[below] {$d/2$} (w3);
    \draw[edge] (w3) -- node[below] {$c/0$} (w4);
    \draw[edge,bend left=18] (w1) to node[right] {$a/0$} (w4);
  \end{scope}
\end{tikzpicture}
\caption{So sánh khái niệm giữa \(F(A)\) và \(F_w(A)\): bản có trọng số lưu định danh bài hát như trọng số trên đường đi thay vì bỏ mất thông tin đầu ra.}
\end{figure}

Đáng chú ý, ngay cả với $15{,}455$ bài hát, tổng số cung của
$F_w(A)$ là $53.0$ triệu, chỉ nhiều hơn khoảng $0.004\%$ so với $F(A)$. Đồng
thời,
\[
  |F(A)|_E \approx 2.1|A|_E.
\]
Quan hệ nhân này được giữ khi kích thước tập bài hát thay đổi từ $1$ tới
$15{,}455$. Hơn nữa, với $15{,}455$ bài hát, $U$ là
$45$-hậu-tố-duy-nhất: số ``va chạm'' hậu tố giảm rất nhanh khi độ dài hậu tố
tăng. Thực nghiệm cũng cho
\[
  |F_w(A)|_Q \approx 28.8\text{ triệu} \approx 1.2|A|_Q,
\]
nghĩa là chặn trong Hệ quả 3 được xác nhận trong bối cảnh thực nghiệm này.

\section{Kết luận}

Bài viết trình bày một phân tích mới về kích thước của ô-tô-mát hậu tố và
ô-tô-mát nhân tố của một tập xâu được biểu diễn bằng ô-tô-mát, theo kích thước
của ô-tô-mát ban đầu. Phân tích cho thấy ô-tô-mát hậu tố và ô-tô-mát nhân tố
có thể thực tế để xây dựng chỉ mục cho số lượng xâu rất lớn. Ứng dụng vào một
bài toán nhận dạng nhạc quy mô lớn cũng minh họa thêm điều này. Ô-tô-mát nhân
tố của ô-tô-mát nhiều khả năng là một chỉ mục hữu ích và gọn cho các tác vụ
quy mô rất lớn.

Bài viết cũng đưa ra thuật toán thời gian tuyến tính để xây dựng ô-tô-mát hậu
tố hoặc ô-tô-mát nhân tố của một tập xâu trong thời gian tuyến tính theo kích
thước của cây tiền tố biểu diễn chúng. Thuật toán áp dụng cho mọi ô-tô-mát đầu
vào hậu-tố-duy-nhất và tổng quát hóa chặt chẽ phép xây dựng trực tuyến chuẩn
của ô-tô-mát hậu tố cho một xâu đơn. Thuật toán và phân tích này đặt ra câu
hỏi tự nhiên về cách xây dựng hiệu quả ô-tô-mát hậu tố của một ô-tô-mát đầu
vào bất kỳ.

\section*{Lời cảm ơn}

Các tác giả cảm ơn Cyril Allauzen vì nhiều trao đổi về nội dung trình bày.
Nghiên cứu của Mehryar Mohri và Eugene Weinstein được hỗ trợ một phần bởi New
York State Office of Science Technology and Academic Research (NYSTAR). Dự án
cũng được tài trợ một phần bởi Department of the Army, số tài trợ
W81XWH-04-1-0307. U.S. Army Medical Research Acquisition Activity, 820 Chandler
Street, Fort Detrick MD 21702-5014 là cơ quan trao và quản lý tài trợ. Nội
dung tài liệu không nhất thiết phản ánh lập trường hay chính sách của Chính
phủ, và không nên suy diễn thành sự xác nhận chính thức.

\begin{thebibliography}{99}
\bibitem{gusfield}
D. Gusfield, \textit{Algorithms on Strings, Trees, and Sequences}, Cambridge
University Press, Cambridge, UK, 1997.

\bibitem{crochemore-rytter}
M. Crochemore, W. Rytter, \textit{Jewels of Stringology}, World Scientific,
2002.

\bibitem{allauzen-index}
C. Allauzen, M. Mohri, M. Saraclar, General Indexation of Weighted Automata:
Application to Spoken Utterance Retrieval, in: Proceedings of the Workshop on
Interdisciplinary Approaches to Speech Indexing and Retrieval (HLT/NAACL),
Boston, Massachusetts, 2004, pp. 33--40.

\bibitem{weinstein-music}
E. Weinstein, P. Moreno, Music Identification with Weighted Finite-State
Transducers, in: Proceedings of the International Conference on Acoustics,
Speech, and Signal Processing (ICASSP), Honolulu, Hawaii, 2007, pp. 689--692.

\bibitem{blumer-smallest}
A. Blumer, J. Blumer, D. Haussler, A. Ehrenfeucht, M. T. Chen, J. I. Seiferas,
The smallest automaton recognizing the subwords of a text, \textit{Theoretical
Computer Science} 40 (1985) 31--55.

\bibitem{crochemore-transducers}
M. Crochemore, Transducers and repetitions, \textit{Theoretical Computer
Science} 45 (1986) 63--86.

\bibitem{blumer-complete}
A. Blumer, J. Blumer, D. Haussler, R. M. McConnell, A. Ehrenfeucht, Complete
inverted files for efficient text retrieval and analysis, \textit{Journal of
the ACM} 34 (1987) 578--589.

\bibitem{inenaga}
S. Inenaga, H. Hoshino, A. Shinohara, M. Takeda, S. Arikawa, G. Mauri,
G. Pavesi, On-line construction of compact directed acyclic word graphs,
\textit{Discrete Applied Mathematics} 146(2) (2005) 156--179.

\bibitem{mohri-robust}
M. Mohri, P. Moreno, E. Weinstein, Robust music identification, detection, and
analysis, in: Proceedings of the International Conference on Music Information
Retrieval (ISMIR), Vienna, Austria, 2007, pp. 135--139.

\bibitem{allauzen-oracle}
C. Allauzen, M. Crochemore, M. Raffinot, Efficient experimental string matching
by weak factor recognition, in: Proceedings of the 12th Annual Symposium on
Combinatorial Pattern Matching (CPM), Springer-Verlag, London, UK, 2001,
pp. 51--72.

\bibitem{revuz}
D. Revuz, Minimisation of acyclic deterministic automata in linear time,
\textit{Theoretical Computer Science} 92 (1992) 181--189.

\bibitem{blumer-thesis}
J. A. Blumer, Algorithms for the directed acyclic word graph and related
structures, Ph.D. thesis, Denver University, 1985.

\bibitem{mohri-transducers}
M. Mohri, Finite-state transducers in language and speech processing,
\textit{Computational Linguistics} 23(2) (1997) 269--311.

\bibitem{mohri-applied}
M. Mohri, Statistical Natural Language Processing, in: M. Lothaire (Ed.),
\textit{Applied Combinatorics on Words}, Cambridge University Press, 2005.
\end{thebibliography}

\end{document}
